\documentclass[10pt]{article}
\usepackage[T1]{fontenc}
\usepackage{lmodern,amsmath,amssymb,amsthm,mathtools,booktabs,microtype}
\usepackage[margin=.92in]{geometry}
\usepackage{xcolor}
\definecolor{ink}{RGB}{25,65,84}
\usepackage[colorlinks=true,linkcolor=ink,urlcolor=ink,citecolor=ink]{hyperref}
\newtheorem{theorem}{Theorem}
\newtheorem{lemma}[theorem]{Lemma}
\newcommand{\R}{\mathbb R}
\newcommand{\Q}{\mathbb Q}
\newcommand{\Z}{\mathbb Z}
\newcommand{\conv}{\operatorname{conv}}
\newcommand{\Res}{\operatorname{Res}}
\setlength{\emergencystretch}{2em}
\title{A transcendental gap for the frequencies $1,4,7$\\[3pt]
\large An exact contact formula for three frequencies in arithmetic progression}
\author{}\date{}
\begin{document}
\maketitle\vspace{-2.3em}
\begin{abstract}
How long can a real trigonometric polynomial with frequencies $\pm1,\pm4,\pm7$ stay strictly positive? In circle units the supremum is $0.618557753484020625\ldots$. Its cosine after multiplication by $\pi$ has algebraic degree $54$, and the gap itself is transcendental. A polynomial with two double contacts and a positive six-point averaging measure certify the exact value. We also give a finite algebraic rule for the gap for every three-term progression $N,N+b,N+2b$. Its candidates come from tangencies of a polynomial curve and two explicit angular grids. This addresses the three-frequency case of a question of Kozma and Oravecz. The proof uses the positive-measure duality developed in the ProveIt article and is independent of hypotheses on zeros of $L$-functions.
\end{abstract}

\section{The gap and its dual}
For a finite nonempty set $S$ of positive integers, write $M(S)$ for the supremum of the lengths of arcs of $\R/\Z$ on which some real trigonometric polynomial with frequencies in $\pm S$ is strictly positive. For $S=\{N+jb:0\le j\le K\}$ with $b<2N$, Kozma and Oravecz proved $M(S)=(K+1)/(2N+Kb)$. Their Question~1 asks for the behavior when $b\ge 2N$ \cite{KO}. The ProveIt article gives an exact two-frequency result, an all-$K$ resonance case and, off resonance, a sharp-order large-$K$ estimate \cite{repo}. We treat $K=2$ for arbitrary positive $N,b$. The answer is a finite rule involving real roots and signs. The example below explains one obstacle to a rational closed formula.

Put $V_S(x)=(T_n(x))_{n\in S}$, where $T_n(\cos t)=\cos(nt)$. After translation to the center of an arc, replacing $f(t)$ by $(f(t)+f(-t))/2$ preserves strict positivity there. The cosine coordinate $x=\cos t$, in radians, maps an arc of circle length $L$ to $[\cos(\pi L),1]$. For $a\in[-1,1]$ write $C_a=\conv V_S([a,1])$.
\begin{lemma}\label{dual}
If $a_*=\max\{a:0\in C_a\}$, then $M(S)=\arccos(a_*)/\pi$. There is a nonzero $h\in\R^{|S|}$ with $h\cdot V_S\ge0$ on $[a_*,1]$. The origin is a convex combination of at most $|S|$ distinct contact vectors in its supporting plane, and every such combination with positive weights uses $a_*$.
\end{lemma}
\begin{proof}
The uniform circle measure shows that $0\in C_{-1}$. The sets are compact and vary continuously in Hausdorff distance with $a$, by uniform continuity of $V_S$. Also $V_S(1)=(1,\ldots,1)$. Strict separation of $0$ from $C_a$ is equivalent to strict positivity of a cosine polynomial on $[a,1]$. Hence the stated maximum exists and gives the gap. If $0$ were interior to $C_{a_*}$, every unit direction would have support function bounded below by a common positive number. This remains true on shrinking the interval slightly, a contradiction. A supporting hyperplane through $0$ now exists. Any representing measure is supported on its contacts. Carath\'eodory's theorem inside that hyperplane supplies the bound on the number of atoms. A finite positive representation using only larger parameters would contradict maximality.
\end{proof}

This duality also gives an arithmetic observation for every integer spectrum. The condition $0\in C_a$ has a real polynomial description using at most $|S|+1$ nodes, nonnegative weights of sum one, and moment equations. Real quantifier elimination makes $a_*$ algebraic. Consequently $M(S)$ is either rational or transcendental. Indeed $z=\exp(i\pi M(S))$ satisfies $z^2-2a_*z+1=0$. If $M(S)$ were algebraic irrational, Gelfond--Schneider applied to $(-1)^{M(S)}$ would make this $z$ transcendental. Semialgebraic computability and fixed-arc positivity tests already follow from general methods, including the interval positivity methods in \cite{Dum}. The following theorem describes the contact configurations explicitly.

\section{A finite rule for three frequencies}
Divide $N,b$ by $g=\gcd(N,b)$ in the construction below and divide the resulting gap by $g$. This follows by pulling a polynomial back under multiplication by $g$ on the circle. Thus take $\gcd(N,b)=1$, write $m=N+b$, and put
\[
 V(x)=(T_N(x),T_m(x),T_{N+2b}(x)),\qquad D(x)=V(x)\times V'(x).
\]
A node set is feasible if the origin is in the convex hull of its images under $V$. The test for at most three nodes requires only linear algebra and sign comparisons. All real roots used below are in $[-1,1]$.

\begin{theorem}\label{algorithm}
Form the following finite list of feasible node sets.
\begin{enumerate}
\item Include every feasible singleton or pair with nodes in the union of the two grids
\[
 \{\cos(j\pi/b):0\le j\le b\},\qquad
 \{\cos((2j+1)\pi/(2m)):0\le j<m\}.
\]
\item Let $E$ consist of $1$ and the common real zeros of the coordinates of $D$, retaining parameters with $V\ne0$. Generate the following nonzero normals, identified up to scaling,
\begin{align*}
 &V(u)\times V(v), &&u,v\in E,\\
 &D(v), &&u\in E,\quad D(v)\cdot V(u)=0,\quad V(u)\times V(v)\ne0,\\
 &D(u), &&D(u)\times D(v)=0,\quad D(u),D(v)\ne0,\quad V(u)\times V(v)\ne0 .
\end{align*}
For each normal $h$, take the distinct real roots of $h\cdot V(x)$. Include the feasible triples among them.
\end{enumerate}
The maximum of the smallest node over all listed sets equals $a_*$. The rule involves univariate root computations and finitely many isolated solutions of bivariate polynomial equations.
\end{theorem}

\begin{proof}
Every listed minimum is feasible as a left endpoint. Choose a positive representation of $0$ on $[a_*,1]$ with the fewest nodes. Lemma~\ref{dual} bounds their number by three.

A singleton is a common zero and belongs to the second grid. For a pair, order the angles $0\le\phi<\theta\le\pi$. Their vectors are antiparallel. If the middle coordinate at $\theta$ vanishes, it vanishes at $\phi$ too, giving two nodes in the second grid. Otherwise
\[
 \cos(Nt)+\cos((N+2b)t)=2\cos(bt)\cos(mt)
\]
forces $\cos(b\theta)=\cos(b\phi)$. If $\sin(b\theta)=0$, both nodes lie in the first grid. In the remaining case write $b\phi\equiv\sigma b\theta\pmod{2\pi}$ with $\sigma=\pm1$. The first two coordinates and $\sin(b\theta)\ne0$ determine both the cosine and sine of the phase $N\phi$, up to the same negative factor as those of $\sigma N\theta$. Since both phase vectors have length one, the factor is $-1$. Thus $d=\phi-\sigma\theta$ has $Nd/\pi$ an odd integer and $bd/\pi$ an even integer. Coprimality makes $d/\pi$ an integer, necessarily odd, with $N$ odd and $b$ even. It follows that $\theta\ge\pi/2$. But $x=0$ is then a feasible singleton, since all three frequencies are odd. This pair cannot give a larger minimum than that listed singleton.

For a minimal triple, no two vectors are proportional. Opposite vectors give a feasible pair, and positively proportional vectors can be combined into one with a rescaling and renormalization of weights. The normal $h$ of the supporting plane therefore depends only on two of the vectors. Each of the two nodes strictly larger than $a_*$ is either the endpoint $1$, a zero of $D$, or an interior regular tangency with $h$ proportional to $D$. This gives exactly the three cases in the displayed list.

The three component polynomials of $V$ have distinct degrees. Thus $h\cdot V$ is nonzero for every constant $h\ne0$. The polynomial vector $D$ cannot vanish identically, since that would make the projective ratios of the components of $V$ constant. For fixed $p=V(u)\ne0$, the polynomial $D(v)\cdot p$ cannot vanish identically. To see this, change linear coordinates so $p=(0,0,1)$. Identical vanishing says the derivative of the ratio of the first two transformed coordinate polynomials is zero, contrary to their linear independence.

Finally suppose a curve of regular pairs $(u,v)$ survived the exclusions in the last line. Work on a smooth piece and choose a smooth normal $h$ by setting one nonzero coordinate equal to one. Differentiating $h\cdot V(u)=h\cdot V(v)=0$ along the piece gives $h'\cdot V(u)=h'\cdot V(v)=0$ because both nodes are tangencies. The vectors $V(u),V(v)$ are independent, hence $h'$ is proportional to $h$ and the normalization gives $h'=0$. A fixed $h$ has finitely many contact parameters, a contradiction. The same argument applies on the smooth locus of a complex algebraic component in characteristic zero. Thus one may divide the three minors of $D(u)\times D(v)$ by their common polynomial gcd to discard their common curve components and intersect the remaining equations, keeping only the regular independent pairs. An excluded component lies in the closed vanishing set of a discarded regularity or independence condition, so this division loses no valid pair. Zero-dimensional polynomial systems have finitely many solutions.
\end{proof}

For example $u=v$ always makes all three minors vanish. Removing this diagonal is essential for isolating the genuine double tangencies. Each normal determined by two inner contacts gives the endpoint by a univariate polynomial equation.

\section{The spectrum $\{1,4,7\}$}
\begin{theorem}\label{example}
Let $M=M(\{1,4,7\})$. Then
\[
 0.61855775348402062 < M < 0.61855775348402063 .
\]
The number $2\cos(\pi M)$ has algebraic degree $54$, with primitive integral minimal polynomial of leading coefficient $729$. In particular $M$ is transcendental. Up to translation, precisely one probability measure of minimum arc length annihilates all six Fourier modes. It gives positive weight to each of six points.
\end{theorem}

For a list of coefficients $[c_0,\ldots,c_d]$ write $[c]_X=c_0X^d+\cdots+c_d$. Let $s$ be the unique root within $10^{-18}$ of $1.12268200790234930920$ of
\[
 H(s)=[165888,-912384,1974528,-2154240,1297800,-467124,
                105672,-15092,1274,-49]_{s^2}.
\]
Define
\begin{align*}
 3924844p={}&[-5579891712,30121984512,-63425084544,66376108800,\\
     &\hspace{18mm}-37602505764,12510048822,-2542346352,303526195,-18986709]_{s^2},\\
 r={}&3s^2-2p-7/4,\qquad e=(-2sr+2s^3+2sp)/(8p^2-1),\\
 q(x)={}&x^2-sx+p,\qquad C(x)=x^3+2sx^2+rx+e.
\end{align*}
Then $M=\arccos(a)/\pi$, where $a$ is the root of $C$ within $10^{-18}$ of $-0.36390801664214987744$. Its other contacts are the roots $u<v$ of $q$.

\begin{proof}
The decisive factorization is
\begin{equation}\label{certificate}
 F(x)=64q(x)^2C(x)=A T_1(x)+B T_4(x)+T_7(x).
\end{equation}
One can verify it by expanding the left side modulo $H(s)$, with $B=64p^2e$ from the constant coefficient and $A=7+[x]F(x)$ from the linear term. The denominator $8p^2-1$ is invertible modulo $H$. Numerically $A=1.6408563\ldots$ and $B=1.0859836\ldots$.
The contact locations and the signs needed for the proof are
\[
\begin{array}{crrr}
\toprule
 \text{Node} & z & T_4(z) & T_7(z)\\
\midrule
 a & -0.363908017 &  0.080867258 &  0.509300244\\
 u &  0.268082359 &  0.466375039 & -0.946360254\\
 v &  0.854599649 & -0.575545813 & -0.777241913\\
\bottomrule
\end{array}
\]
The cubic is strictly increasing on $[a,1]$, since $C'(a)>0$ and $C''(x)=6x+4s>0$ there. Thus $F\ge0$ throughout that interval, with precisely three distinct contacts. In the two projected coordinates $(x,T_4(x))$, the determinants of consecutive columns at $(u,v),(v,a),(a,u)$ lie respectively in
\[
 (-0.554,-0.552),\qquad(-0.141,-0.140),\qquad(-0.192,-0.191).
\]
These determinants give a null vector with all entries of the same sign. Since \eqref{certificate} expresses the third coordinate at the contacts in terms of the first two, this yields positive weights for all three moments after normalization. Every shorter centered arc excludes $a$, and any annihilating measure there would, by integrating $F$, be supported on the two remaining contacts. Their projected determinant is nonzero. Lemma~\ref{dual} proves the exact gap claim. This argument explains why the two interior zeros of $F$ permit strict positivity on every shorter arc.

The three positive weights are uniquely determined by normalization and the two independent moment rows. Their lifts to the centered arc can only be supported on the three pairs of angles $\pm\arccos z$, $z=a,u,v$. The differences between the weights within each pair must satisfy the three homogeneous sine equations. The matrix of these equations, after dividing column $z$ by $\sqrt{1-z^2}$, has entries $U_{n-1}(z)$ for $n=1,4,7$. Here $U_{n-1}(\cos t)=\sin(nt)/\sin t$. Its determinant is positive by the rational bounds checked in the supplement. Thus the weight in each pair splits equally. Translation reduces any arc of the optimum length to this centered one.

For the degree statement, perform arithmetic in $\Q[s]/(H)$ and represent $r,e$ by their polynomial remainders. Take the primitive integral part with positive leading coefficient of
\[
 \Res_s\big(H(s),\,C(y/2)\big)=\text{constant}\cdot Q(y).
\]
The polynomial $Q$ has degree $54$, leading coefficient $729$ and a root at $y=2a$. Its monic reduction $Q_0$ modulo $211$ satisfies
\[
 y^{211^{54}}\equiv y\pmod{Q_0},\qquad
 \gcd(Q_0,y^{211^i}-y)=1\quad(i=18,27).
\]
The prime divisors of $54$ are $2$ and $3$, so these three tests prove irreducibility over $\mathbb F_{211}$ and hence over $\Q$. As $Q$ is primitive and nonmonic, $2a$ is not an algebraic integer. For a rational $M$ one would have $2a=\zeta+\zeta^{-1}$ for a root of unity $\zeta$. Rationality is therefore impossible, and the dichotomy in Section~1 gives transcendence.
\end{proof}

\paragraph{Reproducing the integer and sign checks.}
Run \texttt{python check\_gaps.py} with SymPy. The script isolates $s$ by Sturm counting, verifies the polynomial identity coefficient by coefficient, and brackets the roots of $q,C$ using rational intervals. It computes the three projected signs and the sine determinant and checks the gap enclosure by Taylor bounds with a Machin enclosure for $\pi$. Finally it constructs the degree-$54$ resultant, prints all coefficients and performs the three finite-field checks. No trigonometric approximation enters the proof of optimality.

\paragraph{How to find the certificate.}
For this example the normal $D(t)$ has coordinates
\begin{align*}
 &1536t^{10}-3456t^8+2688t^6-840t^4+112t^2-7,\\
 &-16t^3(24t^4-28t^2+7),\qquad 24t^4-8t^2-1 .
\end{align*}
Equating the first and second ratios to the third coordinate at $u,v$, clearing denominators, and dividing the two antisymmetric minors by $u-v$ gives symmetric polynomial equations. Substitute $s=u+v$, $p=uv$. Elimination gives $H$ among the possible factors in $s$. The linear subresultant determines $p$ on this branch. The two double roots then force the factor $q^2$ in \eqref{certificate}. The sign checks show why this branch supplies the global optimum and allow the other elimination factors to be discarded without computing their real zeros.

\paragraph{Further questions.}
Theorem~\ref{algorithm} covers every finite three-frequency progression in the wide-step range of Question~1. It gives a finite choice of algebraic roots followed by one inverse cosine. The transcendental answer already at $N=1,b=3$ rules out a rational function of the integer parameters valid throughout that range. Which wide progressions nevertheless have rational gaps? What contact patterns and arithmetic degrees occur for four or more frequencies? The double-tangency construction here uses that the supporting space of a three-node rule is a plane.

\begin{thebibliography}{9}
\bibitem{KO} G. Kozma and F. Oravecz,
\emph{On the gaps between zeros of trigonometric polynomials},
Real Analysis Exchange \textbf{28} (2002/03), 447--454.
\href{https://arxiv.org/html/math/0601021v1}{Question 1 and Theorem 3}.
\bibitem{repo} ProveIt Contributors, \emph{Beyond the Riemann hypothesis for Dirichlet $L$-functions},
\href{https://github.com/VladimirReshetnikov/ProveIt/blob/ff96fd44ab160e65fb8a9d861243828d07c156cf/Analysis/GeneralizedRiemannHypothesis/Article/sections/fil.tex\#L1083-L1197}{positive filters, two-frequency gaps and wide arithmetic progressions}, revision \texttt{ff96fd44ab16}.
\bibitem{Dum} B. Dumitrescu, \emph{Positive Trigonometric Polynomials and Signal Processing Applications}, Springer, 2007, Section~1.4.
\href{https://doi.org/10.1007/978-1-4020-5125-8}{DOI and contents}.
\end{thebibliography}
\end{document}
